\documentclass[11pt]{article}
% Detailed primer (for review): Magnus expansions -> Feynman graphs -> a geometric theory of U(3)
% selection for Type-1 N=3 multistate Landau-Zener. Companion to paper/working\_sessions/ws_geom_scope.md and the
% markdown primer paper/drafts/primer_magnus_feynman_geometry.md. Expository; no new claims.
\usepackage[margin=1in]{geometry}
\usepackage{amsmath,amssymb,bm,mathtools}
\usepackage{enumitem}
\usepackage{booktabs}
\usepackage{hyperref}
\usepackage{xcolor}
\hypersetup{colorlinks=true,linkcolor=blue!50!black,urlcolor=blue!50!black}
\newcommand{\dd}{\mathrm{d}}
\newcommand{\ii}{\mathrm{i}}
\newcommand{\eps}{\varepsilon}
\newcommand{\Tord}{\mathcal{T}}
\newcommand{\Wt}{\widetilde W}
\newcommand{\U}{\mathrm{U}}
\newcommand{\Pmid}{P_{m\to m}}
\newcommand{\norm}[1]{\lVert #1 \rVert}
\title{\textbf{From Magnus expansions to a geometric theory of $\U(3)$ selection}\\[2pt]
\large A primer for the Type-1 $N{=}3$ multistate Landau--Zener problem}
\author{Co-theoretical-physicist working note}
\date{June 2026}

\begin{document}
\maketitle
\begin{abstract}
We assemble the conceptual scaffolding behind the WS-GEOM programme: (i) the \emph{Magnus expansion},
which writes a time-ordered propagator as a single exponential of nested commutators, unitary at every
order; (ii) how, in the \emph{adiabatic frame} of the multistate Landau--Zener (MLZ) problem, this becomes
a \emph{Feynman-graph} expansion on the adiabatic sheets --- geometric vertices (the derivative coupling
$W$) joined by dynamical-phase propagators; and (iii) how that structure fleshes out a \emph{geometric
theory} of how the Type-1 Hamiltonian selects its slice of $\U(3)$: a topological permutation/spinor
sector dressed by the $W$-holonomy, bounded by the decoupling locus, with the one irreducible transcendental
$\sigma$ quarantined as the all-orders resummation. The final section describes the four milestones
(M1--M4) of the WS-GEOM workstream. This is an expository companion; it introduces no new claims.
\end{abstract}

\tableofcontents

\section{Why this primer}
The closed-form question for the Type-1 $N{=}3$ middle survival is settled in the negative: the irreducible
content is a single genuinely transcendental Stokes constant $\sigma$, and the most efficient route to the
\emph{exact} value is structured numerical integration of the governing ODE. What remains wide open is a
\emph{geometric} question: the map $\Phi:(\gamma,\eps,a)\mapsto S\in\U(3)$ (the propagator from $u=-\infty$
to $+\infty$) lands the transition data in a particular, structured part of $\U(3)$ --- and we have not
asked \emph{what} that structure is, or \emph{what selects it}. The Magnus/Feynman machinery below is the
natural language for that question, because it separates the rigid (topological, knowable) part of the
answer from the irreducible transcendental part, and it does so $\sigma$-free except for the one resummed
datum we have proven cannot be put in closed form.

\section{Magnus expansions}\label{sec:magnus}

\subsection{The time-ordering problem}
The propagator of a linear system $\dot U(t)=A(t)\,U(t)$, $U(0)=\mathbb 1$, is the \emph{time-ordered
exponential}
\begin{equation}
U(t)=\Tord\exp\!\int_0^t A(s)\,\dd s
=\mathbb 1+\int_0^t\! A(t_1)\,\dd t_1+\int_0^t\!\!\int_0^{t_1}\! A(t_1)A(t_2)\,\dd t_2\,\dd t_1+\cdots,
\label{eq:dyson}
\end{equation}
the \emph{Dyson series}. The obstruction to simply writing $U=\exp\!\int A$ is that $A(t_1)$ and $A(t_2)$
need not commute. Truncating \eqref{eq:dyson} also \emph{breaks unitarity}: a finite sum of ordered
products of anti-Hermitian generators is not in general unitary.

\subsection{The Magnus series}
Magnus's idea \cite{Magnus1954} is to keep a \emph{single} exponential,
\begin{equation}
U(t)=\exp\Omega(t),\qquad \Omega=\Omega_1+\Omega_2+\Omega_3+\cdots,
\end{equation}
where each $\Omega_k$ is built from $k$-fold nested commutators of $A$:
\begin{align}
\Omega_1(t)&=\int_0^t A(t_1)\,\dd t_1,\\
\Omega_2(t)&=\tfrac12\int_0^t\!\!\int_0^{t_1}\big[A(t_1),A(t_2)\big]\,\dd t_2\,\dd t_1,\\
\Omega_3(t)&=\tfrac16\int_0^t\!\!\int_0^{t_1}\!\!\int_0^{t_2}\Big(\big[A(t_1),[A(t_2),A(t_3)]\big]
+\big[A(t_3),[A(t_2),A(t_1)]\big]\Big)\,\dd t_3\,\dd t_2\,\dd t_1.
\end{align}
The series is the continuous analogue of the Baker--Campbell--Hausdorff formula. Two properties make it the
right tool here.

\paragraph{Unitarity at every order.} If $A$ is anti-Hermitian ($A^\dagger=-A$), then every nested
commutator of $A$'s is anti-Hermitian, so each $\Omega_k$ is anti-Hermitian and \emph{every truncation}
$U_{(n)}=\exp(\Omega_1+\cdots+\Omega_n)$ is exactly unitary. This is a built-in consistency check the Dyson
truncation lacks: $S\in\U(3)$ holds order by order, so $P=|S|^2$ stays doubly stochastic automatically.

\paragraph{Convergence and regime.} The series converges for
$\int_0^t\norm{A(s)}\,\dd s<\pi$ \cite{BlanesCasas2009}. Physically: weak coupling or short evolution. In
our application $A=-\ii\Wt$ is the (dynamical-phase-dressed) \emph{adiabatic} coupling, small when the
energy gaps are large --- the \emph{adiabatic} regime (large BE action $\delta$). A Type-1 surprise (M1,
\S\ref{sec:milestones}): the dressed action $\Lambda=\int\norm{\Wt}\approx\pi$ is \emph{marginal almost
everywhere} --- the flow-side face of the permanent-overlap lemma --- so the series is order-improving but
never \emph{fast}, and in the deepest overlap $\Lambda$ exceeds $\pi$ and it ceases to resum (that
resummation is the hard $\sigma$).

\subsection{A two-level sanity check}
For the ordinary two-level LZ problem the leading Magnus term $\Omega_1$ already reproduces the perturbative
(weak-coupling) transition amplitude, and $\Omega_2$ supplies the first phase correction; the full
(non-perturbative) St\"uckelberg result is the resummation. The $N{=}3$ structure below is the multi-sheet
generalization, with the new feature that the sheets can \emph{recombine}.

\section{The adiabatic-$W$ frame and the Feynman-graph structure}\label{sec:feynman}

\subsection{The adiabatic frame}
Write $\psi(u)=\Phi(u)\chi(u)$ with $\Phi=[\phi_1\,\phi_2\,\phi_3]$ the instantaneous eigenframe of
$H(u)=H_0+uA$, $H\Phi=\Phi\,\mathrm{diag}(E_i)$. The Schr\"odinger equation $\ii\psi'=H\psi$ becomes
\begin{equation}
\ii\,\chi'=\big[\,D(u)-\ii\,W(u)\,\big]\chi,\qquad D=\mathrm{diag}(E_i),\quad
W_{ij}=\langle\phi_i|\phi_j'\rangle=\frac{\langle\phi_i|A|\phi_j\rangle}{E_j-E_i}\ (i\neq j),
\label{eq:adiabatic}
\end{equation}
where $W$ is the (anti-Hermitian, off-diagonal) \emph{derivative coupling}. For Type-1 the family shares
one eigenbasis, so \textbf{$W$ is independent of the slope vector $a$} --- it is fixed geometric data.
Passing to the interaction picture with respect to $D$ dresses the coupling by the dynamical phase,
\begin{equation}
\Wt_{ij}(u)=W_{ij}(u)\,\exp\!\Big(\ii\!\int^u\!(E_i-E_j)\,\dd u'\Big),
\end{equation}
and the physical scattering matrix is
$S=\Phi(+\infty)\,\big[\Tord\exp\!\int(-\Wt)\big]\,\Phi(-\infty)^\dagger$ up to diagonal Stark phases that
drop from $P=|S|^2$.

\subsection{Diagrams on three sheets}
Reading the Dyson/Magnus series of $\Tord\exp\!\int(-\Wt)$ as a sum over histories gives a literal
Feynman expansion on the three adiabatic levels:

\begin{center}
\begin{tabular}{ll}
\toprule
\textbf{diagram element} & \textbf{MLZ object}\\
\midrule
line (sheet $i$) & adiabatic level $E_i(u)$ \\
vertex $\Wt_{ij}(u)$ & a non-adiabatic \emph{hop} $j\to i$ at time $u$, amplitude $W_{ij}(u)$ \\
edge (propagator) & accumulated dynamical phase $e^{\ii\int(E_i-E_j)}$ between hops \\
order-$n$ term & sum over all ordered $n$-hop paths across the sheets \\
\midrule
order $0$ (no vertices) & adiabatic following $\Rightarrow$ the \textbf{directed-cycle} permutation $S_0$ \\
$1$ vertex ($\Omega_1$) & one hop $\Rightarrow$ leading LZ/St\"uckelberg amplitude ($\to$ BE factors) \\
$2$ vertices ($\Omega_2$) & two interfering hops $\Rightarrow$ leading $\Pmid$ interference \\
all orders (resummed) & the full holonomy $=\sigma$ (no closed form) \\
\bottomrule
\end{tabular}
\end{center}

\paragraph{Order zero is a permutation.} With $W=0$ the interaction-picture propagator is the identity, so
$S_0=\Phi(+\infty)\,\mathrm{diag(phases)}\,\Phi(-\infty)^\dagger$. Because $\Phi(\pm\infty)$ are the
standard-basis frames \emph{ordered by adiabatic energy}, and that ordering \emph{reverses} between
$u=\mp\infty$, $|S_0|^2$ is a \textbf{permutation matrix} --- the directed $3$-cycle that dominates the
transition matrix at small coupling. This is the geometric skeleton; everything transcendental is in the
$W$-dressing.

\paragraph{The key separation.} Because $W$ is $a$-independent, the \emph{vertices are fixed geometry} and
the \emph{$a$-dependence lives entirely in the phase edges} $e^{\ii\int(E_i-E_j)}$. The Feynman expansion
thus cleanly factors the problem into a geometric part (the vertices, the same for the whole commuting
family) and an elementary kinematic part (the dynamical actions).

\section{Toward a geometric theory of $\U(3)$ selection}\label{sec:geometry}

\subsection{The eigenbundle and its holonomy}
The eigenframe $\Phi(u)$ is a section of the (rank-1$\,\times\,3$) eigenbundle over the $u$-line; $W$ is the
non-abelian Berry/Wilczek--Zee connection on it \cite{Berry1984,WilczekZee1984}, and the scattering matrix
$S$ is its \emph{holonomy} (parallel transport from $-\infty$ to $+\infty$), dressed by the dynamical
phases. ``How does $H$ select its slice of $\U(3)$?'' is therefore ``what is the holonomy representation of
this bundle?''. The Magnus/Feynman structure makes the answer \emph{layered}, and each layer is a geometric
object.

\subsection{Topological sector $\times$ smooth dressing}
\begin{description}[leftmargin=1.4em]
\item[Order $0$ --- a topological datum.] $S_0$ is the permutation by which the adiabatic energy ordering
reverses between $\mp\infty$. That permutation (the directed $3$-cycle) together with the \textbf{spinor
double-cover signs} $\delta_j=\pm1$ --- the sign of each eigenvector, which flips on encircling a branch
point of the spectral curve $\Sigma:\det(E-H(u))=0$ --- is a \emph{discrete, topological label}: the sector
of $\U(3)$ the dynamics lives in, fixed by the geometry of the eigenbundle, not by any dynamical detail.
\item[Higher orders --- the smooth dressing.] The $W$-vertices and phase-edges deform $S$ \emph{within}
that sector. This is the analytic content: the BE factors, the interference, and (resummed) $\sigma$.
\item[The boundary.] As $(\gamma,\eps,a)$ vary, $S$ sweeps a full-dimensional \emph{region} of
reduced-$\U(3)$ (the map is $4\to4$). Its boundary is where a hop amplitude $W_{ij}$ vanishes --- a level
decouples --- which is exactly the classifier locus (elementary $\Leftrightarrow$ a level decouples).
\end{description}

\noindent The geometric theory therefore reads
\begin{equation}
\boxed{\ \U(3)\ \text{selection}\;=\;
\underbrace{\big(\text{topological permutation/spinor sector}\big)}_{\text{eigenbundle monodromy on }\Sigma}
\;\times\;
\underbrace{\big(\text{smooth }W\text{-holonomy dressing}\big)}_{\text{Magnus/Feynman series}},\ }
\end{equation}
with the decoupling locus as the boundary of the swept region, and the hard transcendental $\sigma$
quarantined as the all-orders resummation of the dressing.

\subsection{What a ``geometric invariant'' would be}
Not a closed form for $\sigma$, but a \emph{characteristic-class}-type object: a winding/permutation label,
the parity of the spinor sector, or a relation among the BE actions $\delta_{ij}$ and the turning-point
cross-ratio $\chi$ that is conserved along the family or that cuts the boundary. The Magnus/Feynman
structure tells you \emph{where to look}: order $0$ for the topology; the boundary $W_{ij}\to0$ for the
analytic edge.

\section{The WS-GEOM programme: milestones M1--M4}\label{sec:milestones}
The companion scope (\texttt{paper/working\_sessions/ws\_geom\_scope.md}) organizes the work into four milestones; here is
what each establishes and how it is gated.

\paragraph{M1 --- the adiabatic-$W$ Magnus expansion (\S\ref{sec:feynman}), derivation $+$ numerics.}
Derive $\Omega_0$ (the directed-cycle permutation), $\Omega_1$ (first hop), $\Omega_2$ (first interference)
explicitly, and validate \emph{term by term} against the gold oracle in the \emph{adiabatic} regime.
\emph{Gate [PASSED]:} the error $|P^{(n)}-P^{\mathrm{oracle}}|$ (a) decreases with order $n$ at each
sample, and (b) improves as the system is made more adiabatic (larger $\delta$), up to the marginal
convergence boundary $\Lambda\!\approx\!\pi$; the deepest-overlap non-convergence is confirmed as the honest
boundary. Order $0=$ the directed cycle is coordinator-verified to equal the oracle's dominant pattern and
to be \emph{selected by continuation through the universal node} (energy-sorting gives the wrong
transposition). \emph{Output:} \texttt{experiments/ws\_geom\_magnus.py}, \texttt{paper/working\_sessions/ws\_geom\_m1.md}.

\paragraph{M2 --- map the image region (numerics, facet A).} From large cheap $(\gamma,\eps,a)\to P$ data,
characterize the swept region in reduced-$\U(3)$: the directed-cycle bias (distance to the nearest
permutation vertex vs.\ $\delta$), the effective dimension, and the \emph{boundary} (test that it coincides
with the classifier/decoupling locus). \emph{Gate:} boundary $=$ decoupling locus, quantified.

\paragraph{M3 --- the invariant hunt (caged symbolic regression $+$ derivation).} Search a combination of
$\{\Pmid,b,\delta_{ij},\chi,c_i,\dots\}$ that is conserved along the family or that cuts the boundary.
Symbolic regression is used \emph{only} as a structure detector --- never a surrogate solver, never a
closed-form finder (no closed form exists). Any candidate is handed to \texttt{physics-derivation} to
\emph{prove}. \emph{Gate:} one invariant proven, or precisely refuted (a first-class negative).

\paragraph{M4 --- synthesis (theory).} Assemble M1--M3 into the geometric statement of \S\ref{sec:geometry}:
the section as a topological permutation/spinor sector dressed by the $W$-holonomy, bounded by the
decoupling locus --- written as a paper section, entirely $\sigma$-free.

\paragraph{Honest scope.} The Magnus series does not resum in deep overlap (the same boundary as the
uniform semiclassical law); the geometric \emph{skeleton} (order $0$ $+$ low orders) and the
\emph{topology} are the targets, and they are $\sigma$-free. M3 may return no new invariant --- that too is
a real result (the analytic content of the section is exhausted by BE, double-stochasticity and $\chi$, and
only the topological label remains).

\section{One-paragraph summary}
A Magnus expansion writes the propagator as $e^{\Omega}$ with $\Omega$ a sum of nested commutators of the
coupling --- unitary at every order, convergent in the weak-coupling (generic) regime. For Type-1 MLZ the
coupling is the $a$-independent adiabatic seed $W$, and the expansion becomes a Feynman series on three
sheets: geometric $W$-vertices joined by dynamical-phase edges, summed over hop-histories. Order $0$ is the
directed-cycle permutation set by the eigenbundle's energy-reordering and spinor signs (a \emph{topological}
selection of the $\U(3)$ sector); the higher orders are the smooth $W$-holonomy dressing (the
\emph{analytic} content, with $\sigma$ as its non-closed resummation); the boundary of the swept region is
the decoupling locus. That layered decomposition --- topology $\times$ holonomy dressing, with an explicit
boundary --- is the skeleton of a geometric theory of how the Hamiltonian selects its slice of $\U(3)$.

\begin{thebibliography}{9}
\bibitem{Magnus1954} W.~Magnus, \emph{On the exponential solution of differential equations for a linear
operator}, Comm.\ Pure Appl.\ Math.\ \textbf{7}, 649 (1954).
\bibitem{BlanesCasas2009} S.~Blanes, F.~Casas, J.~A.~Oteo, J.~Ros, \emph{The Magnus expansion and some of
its applications}, Phys.\ Rep.\ \textbf{470}, 151 (2009).
\bibitem{Berry1984} M.~V.~Berry, \emph{Quantal phase factors accompanying adiabatic changes}, Proc.\ R.\
Soc.\ Lond.\ A \textbf{392}, 45 (1984).
\bibitem{WilczekZee1984} F.~Wilczek, A.~Zee, \emph{Appearance of gauge structure in simple dynamical
systems}, Phys.\ Rev.\ Lett.\ \textbf{52}, 2111 (1984).
\end{thebibliography}

\end{document}
